\section{Supplementary Figures}
\label{app:suppfigs}

\subsection{Spider~2.0-Lite Error Taxonomy}
\label{app:error_taxonomy}

Figure~\ref{fig:error_taxonomy} presents the distribution of error types
for zero-shot and few-shot ($k$\,=\,3) configurations on Spider~2.0-Lite.
The primary effect of few-shot retrieval is a reduction in
\texttt{no\_such\_column} errors ($-$19 queries, $-$36\%),
consistent with the hypothesis that few-shot examples guide the LLM
to use the correct column names from the retrieved schema.

\begin{figure*}[!t]
  \centering
  \includegraphics[width=0.68\linewidth, trim=0 8pt 0 8pt, clip]{fig7_error_taxonomy.pdf}
  \caption[Spider~2.0-Lite error taxonomy: zero-shot vs.\ few-shot]{
    Distribution of query outcomes across 135 Spider~2.0-Lite questions
    for zero-shot and few-shot ($k$\,=\,3) configurations.}
  \label{fig:error_taxonomy}
\end{figure*}

\vspace{-4pt}
\subsection{Mixed-Precision Profile Map}
\label{app:mpheatmap}

Figure~\ref{fig:mpheatmap} shows the quantization dtype assigned to each
encoder component under each of the six profiles evaluated.
Token embedding (nn.Embedding) is retained in FP32 under all INT8 profiles
because \texttt{quantize\_dynamic} targets only \texttt{nn.Linear} layers.

\begin{figure*}[!t]
  \centering
  \includegraphics[width=0.85\linewidth]{fig14_mp_heatmap.pdf}
  \caption[Mixed-precision profile map]{
    Per-component quantization dtype (FP32, FP16, INT8) for each of the
    six profiles.
    INT8-Full (right column) is highlighted as the Pareto-optimal choice.}
  \label{fig:mpheatmap}
\end{figure*}

\clearpage
\section{Benchmark Tools}
\label{app:tools}

Table~\ref{tab:tools} lists the benchmark scripts developed for this work.
All scripts are included in the supplementary material.

\begin{table}[!h]
  \centering
  \caption[Benchmark scripts]{Benchmark scripts developed for this work
    and the measurements each produces.}
  \label{tab:tools}
  \begin{tabular}{ll}
    \toprule
    Script & Measurement \\
    \midrule
    \texttt{bench\_latency.py}           & Encode latency + RAM per profile \\
    \texttt{bench\_significance.py}      & Bootstrap CI + McNemar's test \\
    \texttt{bench\_schema\_linking.py}   & R@K, MRR vs.\ quantization profile \\
    \texttt{bench\_layer\_sensitivity.py}& Per-layer INT8 sensitivity \\
    \texttt{bench\_spider2\_rk.py}       & Spider~2.0 few-shot R@K \\
    \bottomrule
  \end{tabular}
\end{table}

\section{Reproducibility}
\label{app:repro}

\textbf{Spider~1.0 best configuration:}
\begin{lstlisting}[language=bash,basicstyle=\small\ttfamily,breaklines=true]
cd spider1_pipeline
EMBED_INT8=true python3 run.py run \
  --split dev --sc_n 3 --k 5 \
  --no_pruning --two-pass --evaluate
\end{lstlisting}

\textbf{Spider~2.0-Lite full pipeline:}
\begin{lstlisting}[language=bash,basicstyle=\small\ttfamily,breaklines=true]
cd spider2/spider2_pipeline
python3 run.py run \
  --few-shot-k 3 --sc 3 --passes 2 --evaluate
\end{lstlisting}

\textbf{Quantization benchmark:}
\begin{lstlisting}[language=bash,basicstyle=\small\ttfamily,breaklines=true]
cd spider1_pipeline
python3 bench_schema_linking.py --device cpu
python3 bench_layer_sensitivity.py --fine
\end{lstlisting}

\section{Hardware Configuration}
\label{app:hardware}

\noindent
\begin{tabular}{ll}
  \toprule
  Component & Specification \\
  \midrule
  Local CPU/MPS & Apple M-series (ARM64, unified memory) \\
  Server GPU    & NVIDIA RTX~3090, 24\,GB VRAM \\
  LLM serving   & Ollama v0.x, \texttt{localhost:11434} \\
  LLM model     & \texttt{qwen3.5:9b} (GGUF, Q4\_K\_M) \\
  Embedding     & arctic-embed-m, 109.5\,M parameters \\
  Python        & 3.12.0 \\
  PyTorch       & 2.10.0 \\
  \bottomrule
\end{tabular}
